\clearpage
\begin{examplebox}
  \textbf{\textcolor{CExample}{例1：先不用公式，直接看哪两点配对（入门）}}
\begin{center}\begin{tikzpicture}[font=\small,>=stealth,line join=round]
\draw[gray!40,rounded corners=3pt] (0,0) rectangle (14.2,8.3);\node at (7.1,7.700000000000001) {左右等距、函数值同高：对称轴就在中间};
\begin{scope}[shift={(1.7923076923076926,2.038732394366197)},x=1.153846153846154cm,y=0.647887323943662cm]\draw[->,gray] (-0.6,0)--(4.6,0) node[right]{$x$};\draw[->,gray] (0,-0.6)--(0,6.5) node[above]{$y$};\draw[blue,thick,solid,domain=-0.1:4.1,samples=180] plot(\x,{(\x-2)^2+1});\fill[orange!85!black] (0,5) circle(2.2pt);\fill[orange!85!black] (1,2) circle(2.2pt);\fill[orange!85!black] (2,1) circle(2.2pt);\fill[orange!85!black] (3,2) circle(2.2pt);\fill[orange!85!black] (4,5) circle(2.2pt);\draw[violet,dashed] (2,-.5)--(2,6);\node[above,violet] at (2,6){$x=2$};\draw[orange!85!black,dashed] (0,5)--(4,5);\draw[orange!85!black,dashed] (1,2)--(3,2);\foreach \v in {0,1,2,3,4}{\draw[gray] (\v,0)--(\v,-.1);\node[below=5pt] at (\v,-.1){$\v$};}\end{scope}\node[anchor=west,align=left,text width=5.0cm] at (8.4,4.0) {$f=(x-2)^2+1$\\[9pt]$f(0)=f(4)=5$\\[9pt]$f(1)=f(3)=2$\\[9pt]每对横坐标之和为 $4$\\[9pt]轴的位置：$4\div2=2$};
\node[text width=12.6cm,align=center] at (7.1,.65) {图示具体配对；完整证明要求定义域左右成对，且每一对都同高。};
\end{tikzpicture}\end{center}


  已知 $f(x)=(x-2)^2+1$，求 $f(0),f(1),f(2),f(3),f(4)$，并判断对称轴。

  \textbf{解：}依次得到 $5,2,1,2,5$。把数据排成一行：
  \[
    \begin{array}{c|ccccc}
      x&0&1&2&3&4\\ \hline
      f(x)&5&2&1&2&5
    \end{array}
  \]
  $0$ 和 $4$ 在 $2$ 的左右各相差 $2$，高度都是 $5$；
  $1$ 和 $3$ 在 $2$ 的左右各相差 $1$，高度都是 $2$。
  但只检查这两对还不够。再令 $t$ 是任意实数：
  \[
    f(2+t)=t^2+1=f(2-t).
  \]
  定义域是全体实数，左右没有缺口，所以对称轴确实是 $x=2$。

  \medskip
  
\end{examplebox}
\clearpage
\begin{examplebox}
\textbf{\textcolor{CExample}{例2：最简单的绝对值函数（巩固）}}
\begin{center}\begin{tikzpicture}[font=\small,>=stealth,line join=round]
\draw[gray!40,rounded corners=3pt] (0,0) rectangle (14.2,8.3);\node at (7.1,7.700000000000001) {例2：尖点在轴上，两侧绝对值相同};
\begin{scope}[shift={(2.3857142857142857,2.161111111111111)},x=0.8571428571428571cm,y=1.0222222222222221cm]\draw[->,gray] (-1.5,0)--(5.5,0) node[right]{$x$};\draw[->,gray] (0,-0.5)--(0,4) node[above]{$y$};\draw[blue,thick,solid,domain=-1:5,samples=180] plot(\x,{abs(\x-2)});\fill[orange!85!black] (1,1) circle(2.2pt);\fill[orange!85!black] (2,0) circle(2.2pt);\fill[orange!85!black] (3,1) circle(2.2pt);\draw[violet,dashed] (2,-.5)--(2,3.5);\node[above,violet] at (2,3.5){$x=2$};\end{scope}\node[anchor=west,align=left,text width=5.0cm] at (8.4,4.0) {$f=|x-2|$\\[9pt]顶点 $(2,0)$\\[9pt]$f(2-t)=|-t|=|t|$\\[9pt]$f(2+t)=|t|$};
\node[text width=12.6cm,align=center] at (7.1,.65) {轴对称不要求曲线光滑；V形两侧在所有等距位置同高。};
\end{tikzpicture}\end{center}


  判断 $f(x)=|x-2|$ 是否关于 $x=2$ 对称。

  \textbf{解：}定义域为全体实数。把 $2+t$ 和 $2-t$ 分别代进去：
  \[
    f(2+t)=|t|,\qquad f(2-t)=|-t|=|t|.
  \]
  每一对位置的函数值相同，因此图象关于 $x=2$ 对称。
  \textbf{直观理解：}这是一个以 $(2,0)$ 为尖点的 V 形图象，两边像镜子一样。

  \medskip
  
\end{examplebox}
\clearpage
\begin{examplebox}
\textbf{\textcolor{CExample}{例3：不画图，直接用“横坐标之和”求值（最快）}}
\begin{center}\begin{tikzpicture}[font=\small,>=stealth,line join=round]
\draw[gray!40,rounded corners=3pt] (0,0) rectangle (14.2,8.5);\node at (7.1,7.9) {例3：两输入的平均数是3，所以高度相同};
\draw[->,gray] (1,2)--(13,2) node[right]{$x$};\draw[violet,dashed] (7,1.5)--(7,6.7);\draw[blue,thick] (2,5)--(12,5);\fill[orange!85!black] (2,5) circle(2.4pt);\draw[dashed,gray] (2,2)--(2,5);\node[below=7pt] at (2,2){$-2$};\fill[orange!85!black] (12,5) circle(2.4pt);\draw[dashed,gray] (12,2)--(12,5);\node[below=7pt] at (12,2){$8$};\node[above,violet] at (7,6.7){$x=3$};\node[above=8pt] at (10,5){$y=7$};
\node[text width=12.6cm,align=center] at (7.1,.65) {只画题目保证的两个点及配对线；不臆造整条未知曲线。};
\end{tikzpicture}\end{center}


  已知函数图象关于直线 $x=3$ 对称，且 $f(8)=7$，求 $f(-2)$。

  \textbf{解：}因为 $8+(-2)=6=2\times3$，所以 $8$ 与 $-2$ 关于 $3$ 配对。
  两点高度相同，直接得到
  \[
    \boxed{f(-2)=f(8)=7}.
  \]
  \textbf{眼前一亮：}不需要知道 $f(x)$ 的具体表达式，就能求出函数值。

  \medskip
  
\end{examplebox}
\clearpage
\begin{examplebox}
\textbf{\textcolor{CExample}{例4：已知对称轴求参数（高考常见）}}
\begin{center}\begin{tikzpicture}[font=\small,>=stealth,line join=round]
\draw[gray!40,rounded corners=3pt] (0,0) rectangle (14.2,8.3);\node at (7.1,7.700000000000001) {例4：参数确定后，横坐标配对的中点为1};
\begin{scope}[shift={(2.9461538461538463,2.068181818181818)},x=1.1538461538461537cm,y=0.6969696969696969cm]\draw[->,gray] (-1.6,0)--(3.6,0) node[right]{$x$};\draw[->,gray] (0,-0.6)--(0,6) node[above]{$y$};\draw[blue,thick,solid,domain=-1.2:3.2,samples=180] plot(\x,{(\x-1)^2});\fill[orange!85!black] (0,1) circle(2.2pt);\fill[orange!85!black] (1,0) circle(2.2pt);\fill[orange!85!black] (2,1) circle(2.2pt);\draw[violet,dashed] (1,-.5)--(1,5.5);\node[above,violet] at (1,5.5){$x=1$};\end{scope}\node[anchor=west,align=left,text width=5.0cm] at (8.4,4.0) {$m=-2$ 时\\[9pt]$f=x^2-2x+1=(x-1)^2$\\[9pt]$f(0)=f(2)=1$\\[9pt]中点横坐标为 $1$};
\node[text width=12.6cm,align=center] at (7.1,.65) {两种解法都要得到同一条轴：展开比较系数，或直接用二次函数轴公式。};
\end{tikzpicture}\end{center}


  已知 $f(x)=x^2+mx+1$ 的图象关于 $x=1$ 对称，求 $m$。

  \textbf{解法一（按本结论来）：}图象关于 $x=1$ 对称，所以
  \[
    f(1+t)=f(1-t)\quad\text{对任意实数 }t\text{ 成立}.
  \]
  两式作差，原来相同的常数项和 $t^2$ 消去：
  \[
    \begin{aligned}
    f(1+t)-f(1-t)
      &=\bigl[(1+t)^2+m(1+t)+1\bigr]\\
      &\quad-\bigl[(1-t)^2+m(1-t)+1\bigr]\\
      &=(4+2m)t.
    \end{aligned}
  \]
  因为不管 $t$ 取多少，差都应是 $0$，取 $t=1$ 即得
  $4+2m=0$，所以 $\boxed{m=-2}$。

  \textbf{解法二（联系二次函数公式）：}对称轴 $x=-\dfrac{m}{2}$，
  所以 $-\dfrac m2=1$，得到 $m=-2$。

  \textbf{验算：}此时 $f(x)=(x-1)^2$，显然关于 $x=1$ 对称。

  \medskip
  
\end{examplebox}
\clearpage
\begin{examplebox}
\textbf{\textcolor{CExample}{例5：只知道右半边，怎样求左半边（综合）}}
\begin{center}\begin{tikzpicture}[font=\small,>=stealth,line join=round]
\draw[gray!40,rounded corners=3pt] (0,0) rectangle (14.2,8.3);\node at (7.1,7.700000000000001) {例5：把已知的右半边，关于竖线配到左侧};
\begin{scope}[shift={(1.7000000000000002,1.8921052631578945)},x=1.2cm,y=0.2421052631578947cm]\draw[->,gray] (-0.5,0)--(4.5,0) node[right]{$x$};\draw[->,gray] (0,-1)--(0,18) node[above]{$y$};\draw[violet,thick,dashed,domain=0:2,samples=180] plot(\x,{pow(2,4-\x)});\draw[blue,thick,solid,domain=2:4,samples=180] plot(\x,{pow(2,\x)});\fill[orange!85!black] (0,16) circle(2.2pt);\fill[orange!85!black] (1,8) circle(2.2pt);\fill[orange!85!black] (2,4) circle(2.2pt);\fill[orange!85!black] (3,8) circle(2.2pt);\fill[orange!85!black] (4,16) circle(2.2pt);\draw[gray,dashed] (2,-.5)--(2,17);\end{scope}\node[anchor=west,align=left,text width=5.0cm] at (8.4,4.0) {蓝：$x\ge2,\ f(x)=2^x$\\[9pt]紫：$x<2,\ f(x)=2^{4-x}$\\[9pt]$0\leftrightarrow4$，同为 $16$\\[9pt]$1\leftrightarrow3$，同为 $8$\\[9pt]$4-x>2$ 才能用右式};
\node[text width=12.6cm,align=center] at (7.1,.65) {图示区间 $[0,4]$；完整函数定义在实数上，左侧指数中的自变量是 $4-x$。};
\end{tikzpicture}\end{center}


  已知 $f(x)$ 的图象关于直线 $x=2$ 对称，且当 $x\ge2$ 时
  $f(x)=2^x$。求 $x<2$ 时的表达式。

  \textbf{第一步：}左侧任取 $x<2$，镜子另一侧的横坐标是 $4-x$。
  因为 $x<2$，所以 $4-x>2$，恰好进入已知的右半边。

  \textbf{第二步：}图象关于 $x=2$ 对称，故
  \[
    f(x)=f(4-x)=2^{4-x}\qquad(x<2).
  \]
  \textbf{第三步：}例如 $x=0$ 对应 $x=4$，
  $f(0)=2^4=16=f(4)$。$x=1$ 对应 $x=3$，
  $f(1)=2^3=8=f(3)$。

  \textbf{答案：}
  \[
     \boxed{f(x)=\begin{cases}2^{4-x},&x<2,\\2^x,&x\ge 2.\end{cases}}
  \]
  \textbf{方法一句话：}把未知半边的 $x$ 换成 $2a-x$，就回到已知半边。

  \medskip
  
\end{examplebox}
\clearpage
\begin{examplebox}
\textbf{\textcolor{CExample}{例6：给出的式子不是标准写法，怎么直接读出对称轴？}}
\begin{center}\begin{tikzpicture}[font=\small,>=stealth,line join=round]
\draw[gray!40,rounded corners=3pt] (0,0) rectangle (14.2,8.5);\node at (7.1,7.9) {例6：真实代入位置相加恒为6，中点才是轴};
\draw[->,gray] (1,2)--(13,2) node[right]{$x$};\draw[violet,dashed] (7,1.5)--(7,6.7);\draw[blue,thick] (3,4)--(11,4);\fill[orange!85!black] (3,4) circle(2.4pt);\draw[dashed,gray] (3,2)--(3,4);\node[below=7pt] at (3,2){$t+1$};\fill[orange!85!black] (11,4) circle(2.4pt);\draw[dashed,gray] (11,2)--(11,4);\node[below=7pt] at (11,2){$5-t$};\node[above,violet] at (7,6.7){$x=3$};\node[above=8pt] at (10,4){同高};
\node[text width=12.6cm,align=center] at (7.1,.65) {横向以 $t=0$ 的配对 $1,5$ 示意；一般位置之和恒为6，不指定函数值。};
\end{tikzpicture}\end{center}


  已知某函数的定义域为全体实数，且对所有实数 $t$ 有
  \[
      f(t+1)=f(5-t).
  \]
  求它的一条对称轴。

  \textbf{解：}左边位置是 $t+1$，右边位置是 $5-t$；
  它们的和恒为 $6$，平均数恒为 $3$。
  两边横坐标成对、函数值又相等，故\textbf{对称轴是 $x=3$}。
  也可令 $u=t+1$，则 $5-t=6-u$，等式变为
  \[
      f(u)=f(6-u).
  \]
  \textbf{提醒：}这里\textbf{不是}看 $t$ 的系数来猜对称轴，而是先找实际代入 $f$ 的两个横坐标。

  \medskip
  
\end{examplebox}
\clearpage
\begin{examplebox}
\textbf{\textcolor{CExample}{例7：两条对称轴，居然能推出周期（提高）}}
\begin{center}\begin{tikzpicture}[font=\small,>=stealth,line join=round]
\draw[gray!40,rounded corners=3pt] (0,0) rectangle (14.2,8.3);\node at (7.1,7.700000000000001) {例7：两条轴相距3，反射两次横移6};
\begin{scope}[shift={(2.242857142857143,3.7966666666666664)},x=0.5714285714285714cm,y=1.5333333333333332cm]\draw[->,gray] (-2,0)--(8.5,0) node[right]{$x$};\draw[->,gray] (0,-1.4)--(0,1.6) node[above]{$y$};\draw[blue,thick,solid,domain=-1.5:8,samples=180] plot(\x,{cos(deg(pi*(\x-1)/3))});\fill[orange!85!black] (0,0.5) circle(2.2pt);\fill[orange!85!black] (2,0.5) circle(2.2pt);\fill[orange!85!black] (6,0.5) circle(2.2pt);\foreach \a in {1,4}{\draw[violet,dashed] (\a,-1.3)--(\a,1.3);\node[above,violet] at (\a,1.3){$x=\a$};}\end{scope}\node[anchor=west,align=left,text width=5.0cm] at (8.4,4.0) {示例 $f=\cos\frac{\pi(x-1)}3$\\[9pt]轴 $x=1,\ x=4$\\[9pt]配对 $0\to2\to6$\\[9pt]三处高度都为 $\frac12$\\[9pt]$f(x+6)=f(x)$};
\node[text width=12.6cm,align=center] at (7.1,.65) {这条示例曲线帮助看反射；正文代数推导对所有满足题设的函数成立。};
\end{tikzpicture}\end{center}


  已知 $f(x)$ 的定义域为全体实数，图象同时关于 $x=1$、$x=4$ 对称。
  证明 $f(x+6)=f(x)$。

  \textbf{第一步：}关于 $x=1$ 对称，意味着
  \[
      f(x)=f(2-x).
  \]
  \textbf{第二步：}关于 $x=4$ 对称，意味着任意 $u$ 都有
  $f(u)=f(8-u)$。令 $u=2-x$，则
  \[
      f(2-x)=f\bigl(8-(2-x)\bigr)=f(x+6).
  \]
  \textbf{合起来：}
  \[
      \boxed{f(x+6)=f(x)}.
  \]
  两条对称轴相距 $3$，左右翻折两次等于平移 $6$，所以 $6$ 是周期。
  \textbf{注意：}只能肯定 $6$ 是一个周期，不能仅凭题目就说它是\textbf{最小正周期}。
\end{examplebox}
